\section*{局限性}
\phantomsection
\label{sec:limitations}

本工作研究的是一个前景可观但方差很大的设定，因此我们结论的适用范围应当据此加以理解。


\paragraph{基准范围。}
我们的评测在 Terminal-Bench~2 上驱动演化，并在 SWE-bench-verified 上考察迁移能力。尽管冻结后的 harness 能够迁移到第二个任务面以及另外三个基座模型家族，但更广泛的编程语言、仓库规模的部署，以及人在回路的工作流，仍未得到检验。

\paragraph{演化工作点。}
AHE 的步数预算与单任务超时是在演化过程中针对 GPT-5.4 high 拟合的，因此跨模型迁移的数字把 harness 的可移植性与工作点耦合混在了一起——在同一家族内部，增益在不同推理档位之间是非单调的（\S\ref{sec:experiments:transfer}）。要厘清这两个因素，需要在多个工作点下重新运行整个闭环。


\paragraph{自修改治理。}
AHE 把编辑限定在一个工作区内，把每一处改动都记录到带版本的变更清单中进行归因，并以文件粒度回滚无效的编辑，但它并未提供一套完整的护栏体系。长程的 harness 清理与更强的滥用防范仍不完善，因此 AHE 应被视为一个受控的研究原型，而非一个完全成熟的自主自我改进系统。
